\section{Proofs for Section~\ref{sec:identification}}\label{app:idproofs}

Throughout, a stationary $M$-variate point process $\bm N$ with mean
$\bar{\bm\Lambda}$ has a covariance measure
$\operatorname{diag}(\bar{\bm\Lambda})\delta_0(\dd u)+\bm c(u)\,\dd u$, and Bartlett
spectrum \eqref{eq:th-bartlett-def}. All Bartlett spectra below have continuous
densities. Equality of spectra therefore means equality of these densities at every
$\omega$.

\subsection{Auxiliary lemmas}\label{app:th-lemmas}

\begin{lemma}[Cox processes and independent thinning]\label{lem:th-cox2}
\begin{enumerate}[(a)]
  \item Let $\bm N$ be an $M$-variate stationary Cox process directed by
  $\bm\xi=(\xi_k)$. Assume $\E\xi_k(0)^2<\infty$, write $\bar{\bm\Lambda}$ for the mean
  and $\bm c_\xi(u)=\Cov\{\bm\xi(t+u),\bm\xi(t)\}$ for the covariance, and assume
  $\bm c_\xi$ has integrable entries. Then $\bm N$ has covariance density $\bm c_\xi$
  and Bartlett spectrum
  $(2\pi)^{-1}[\operatorname{diag}(\bar{\bm\Lambda})+\hat{\bm c}_\xi(\omega)]$, where
  $\hat{\bm c}_\xi(\omega)$ is Hermitian and positive semidefinite for every $\omega$.
  \item Let $\bm N$ be stationary with mean $\bar{\bm\Lambda}$ and covariance density
  $\bm c$. Record each point of stream $k$ independently with probability $p_k$, and
  let $\bm P=\operatorname{diag}(p_k)$. The recorded process $\bm N^{\bm P}$ has mean
  $\bm P\bar{\bm\Lambda}$ and covariance density $\bm P\bm c(u)\bm P$.
\end{enumerate}
\end{lemma}

\begin{proof}
(a) Given $\bm\xi$, the streams are independent Poisson processes. Hence, for bounded
Borel sets $I,J$,
\[
\E[N_k(I)N_j(J)\mid\bm\xi]=\delta_{kj}\int_{I\cap J}\xi_k
+\int_I\xi_k\int_J\xi_j .
\]
Taking expectations gives
$\Cov\{N_k(I),N_j(J)\}=\delta_{kj}\bar\Lambda_k|I\cap J|+\int_I\int_J c_{\xi,kj}(t-s)\,\dd t\,\dd s$.
For $\bm w\in\mathbb C^M$, the function $\bm w^*\bm c_\xi(u)\bm w$ is the covariance
function of the stationary process $\bm w^*\bm\xi$. It is integrable, so by Fej\'er's
theorem its Fourier transform is
\[
\lim_{T\to\infty}T^{-1}\,\E\Big|\int_0^Te^{-i\omega t}\{\bm w^*\bm\xi(t)-\E\bm w^*\bm\xi(t)\}\,\dd t\Big|^2\ge0 .
\]

(b) Let $R_i$ be the recording indicators. Pairs of distinct points contribute
$\E[R_iR_{i'}]=p_kp_j$ times the second factorial moment measure, whose density off the
diagonal is $\bar\Lambda_k\bar\Lambda_j+c_{kj}$. Each point paired with itself
contributes $p_k\bar\Lambda_k|I\cap J|$ when $k=j$. Subtracting
$p_kp_j\bar\Lambda_k\bar\Lambda_j|I||J|$ gives the claim.
\end{proof}

\begin{lemma}[Poisson cluster processes]\label{lem:th-cluster}
Let $N$ be a stationary Poisson cluster process on $\mathbb R$. Cluster centres form a
Poisson process of rate $r$. Each centre $x$ carries an independent copy of a finite
random configuration $\{y_1,\dots,y_{\mathcal S}\}$, possibly empty, placed at $x+y_i$. Assume
$\E e^{u_0\mathcal S}<\infty$ for some $u_0>0$. Then:
\begin{enumerate}[(a)]
  \item for every $j\ge1$,
  $\chi_j=\lim_{T\to\infty}T^{-1}\operatorname{cum}_jN(0,T]=r\,\E \mathcal S^j$;
  \item the law of $N$ determines $w(k)=r\Pr(\mathcal S=k)$ for every $k\ge1$.
\end{enumerate}
\end{lemma}

\begin{proof}
(a) Let $Y(x)$ be the number of points of a cluster centred at $x$ that fall in
$(0,T]$. The marked Poisson process of centres has Laplace functional
\[
\log\E e^{-sN(0,T]}=-r\int_{\mathbb R}\E[1-e^{-sY(x)}]\,\dd x,\qquad s\ge0 .
\]
For every cluster, $Y(x)^j\le \mathcal S^{j-1}Y(x)$ and $\int Y(x)\,\dd x=\mathcal S\,T$. Hence
$\int\E Y(x)^j\,\dd x\le T\,\E \mathcal S^j<\infty$. By dominated convergence, the right-hand
side is $j$ times differentiable at $s=0+$, with $j$th derivative
$(-1)^jr\int\E Y(x)^j\,\dd x$. Therefore $\E N(0,T]^j<\infty$ by monotone convergence,
and
\[
\operatorname{cum}_jN(0,T]=r\int_{\mathbb R}\E Y(x)^j\,\dd x .
\]
Let $\operatorname{diam}=\max_iy_i-\min_iy_i$. Then $Y(x)=\mathcal S$ on an interval of
$x$-values of length $(T-\operatorname{diam})_+$, so
\[
\mathcal S^j(1-\operatorname{diam}/T)_+\le T^{-1}\int Y(x)^j\,\dd x\le \mathcal S^j .
\]
Dominated convergence gives $T^{-1}\operatorname{cum}_jN(0,T]\to r\E \mathcal S^j$.

(b) Suppose two such processes have the same law. By (a), $r\E \mathcal S^j=r'\E \mathcal S'^j$ for every
$j\ge1$. The finite measures $\nu(k)=k\,w(k)$ and $\nu'(k)=k\,w'(k)$ on
$\{1,2,\dots\}$ then have the same moments of every order $j\ge0$. Both have finite
exponential moments for $u<u_0$. Their generating functions
$\sum_k\nu(k)z^k$ are therefore analytic for $|z|<e^{u_0}$ and agree, so
$\nu=\nu'$ and hence $w=w'$.
\end{proof}

By \cref{prop:stationary}, a stationary Hawkes process with any kernel satisfies
\cref{lem:th-cluster} with $r=\mu$ and $\mathcal S\sim$ Borel$(n)$. The process
$C(\nu,r,a,\gamma)$ superposes a Poisson process of rate $\nu$ (clusters of size $1$)
and a Neyman--Scott process with centre rate $r$ and $\mathcal S\sim\Pois(m)$. It therefore satisfies
the lemma with $w_C(k)=\nu\ind{k=1}+r e^{-m}m^k/k!$. Independent $p$-thinning of a
Poisson cluster process gives a Poisson cluster process with cluster size
$\Bin(\mathcal S,p)$.

\begin{lemma}[Janossy densities identify the intensity]\label{lem:th-janossy}
Let $\bm N$ be an $M$-variate point process on $[0,T]$ with $\E\bm N[0,T]<\infty$. Its
intensity, with respect to its own history and a deterministic calendar, is
$\lambda_k(t)=h_k(t;\bm N|_{[0,t)})$ for measurable functionals with
$h_k\ge h_{\min}>0$. For a configuration $x$ on $[0,t]$ with points
$(s_1,k_1),\dots,(s_m,k_m)$, $s_1<\dots<s_m$, the Janossy density of
$\bm N|_{[0,t]}$ is
\[
j_t(x)=\prod_{i=1}^m h_{k_i}(s_i;x_{<s_i})\exp\Big\{-\sum_k\int_0^t h_k(u;x_{<u})\,\dd u\Big\}>0
\]
\citep[Ch.~7]{daley2003}. For a configuration $x$ on $[0,t)$,
$h_k(t;x)=j_t(x\cup\{(t,k)\})/j_t(x)$. Consequently, if two such specifications
$h^{(1)},h^{(2)}$ give the same law on $[0,T]$, then
$h^{(1)}_k(t;x)=h^{(2)}_k(t;x)$ for almost every $(t,x)$ with respect to Lebesgue
measure times the Lebesgue--Poisson measure. In particular this holds for almost every
$t$ at $x=\emptyset$, and for almost every $(s,t)$ at $x=\{(s,j)\}$.
\end{lemma}

\begin{proof}
Adding a point at the right end $t$ multiplies $j_t$ by $h_k(t;x)$ and leaves the
integral unchanged, which gives the ratio formula. The law determines its Janossy
densities up to null sets. Since $j_t>0$, almost-everywhere statements under the law
transfer to the Lebesgue--Poisson reference measure. The empty configuration is an atom
of that measure, and the one-point configurations carry the measure $\dd s$.
\end{proof}

\begin{lemma}[Affine observable intensity characterises the Hawkes law]\label{lem:th-affine}
Let $H=H(\mu,\varphi)$ be a stationary Hawkes process with kernel
$\varphi\in L^1\cap L^2$, $\int\varphi=n<1$. Let $N$ be a stationary point process,
second-order equivalent to $H$, whose intensity with respect to its own history is
$\lambda(t)=c+\int_{(-\infty,t)}h(t-s)\,\dN(s)$, with $c>0$, $h\ge0$,
$h\in L^1\cap L^2$ and $\int h<1$. Then $h=\varphi$ almost everywhere, $c=\mu$, and
$\Pr_N=\Pr_H$.
\end{lemma}

\begin{proof}
Let $c_N$ be the covariance density of $N$. Since $\dN(t)-\lambda(t)\dd t$ is
orthogonal to $\Hist_{t-}$, we have $\Cov\{\dN(t)-\lambda(t)\dd t,\dN(s)\}=0$ for
$s<t$. Expanding $\Cov\{\lambda(t),\dN(s)\}$ gives, for $u>0$,
\begin{equation}\label{eq:th-wh}
c_N(u)=\bar\Lambda h(u)+\int_0^\infty h(v)\,c_N(u-v)\,\dd v .
\end{equation}
The kernel $\varphi$ solves \eqref{eq:th-wh} with the same $c_N=c_H$, since $H$ has
affine intensity with kernel $\varphi$. The difference $d=h-\varphi\in L^1\cap L^2$ is
supported on $(0,\infty)$ and satisfies
$\bar\Lambda d(u)+\int_0^\infty d(v)c_N(u-v)\,\dd v=0$ for $u>0$. Multiply by $d(u)$,
integrate over $u>0$ and apply Parseval:
\[
0=\frac1{2\pi}\int|\hat d(\omega)|^2\{\bar\Lambda+\hat c_N(\omega)\}\,\dd\omega
=\int|\hat d(\omega)|^2f_H(\omega)\,\dd\omega .
\]
Since $|1-\hat\varphi|\le1+n$, we have $f_H\ge\bar\Lambda(1+n)^{-2}/(2\pi)>0$. Hence
$\hat d=0$, and so $d=0$. Taking expectations of the intensity gives
$c=\bar\Lambda(1-n)=\mu$. Thus $N$ is a stationary solution of the Hawkes equation with
parameters $(\mu,\varphi)$, and $\Pr_N=\Pr_H$ by the uniqueness result of
\citet{bremaud1996}.
\end{proof}

\subsection{Proof of \texorpdfstring{\cref{thm:equiv}}{Theorem (equivalence)}}

(a) By \eqref{eq:th-fH} and \eqref{eq:th-fC}, the means are $\bar\Lambda$ and
$\nu+r a/\gamma$. Given equal means, the spectra agree if and only if
\[
\frac{\bar\Lambda\beta^2n(2-n)}{\beta^2(1-n)^2+x}=\frac{r a^2}{\gamma^2+x}
\qquad\text{for all }x=\omega^2\ge0 .
\]
Both numerators are positive. Cross-multiplying and comparing coefficients in $x$
gives $r a^2=\bar\Lambda\beta^2n(2-n)$ and $\gamma^2=\beta^2(1-n)^2$. The converse is
immediate.

(b) Solving \eqref{eq:th-match} for given $a$ gives
$r=\bar\Lambda\beta^2n(2-n)/a^2$ and
\[
\nu=\bar\Lambda-\frac{r a}{\gamma}
=\bar\Lambda\Big(1-\frac{\beta n(2-n)}{a(1-n)}\Big),
\]
which is nonnegative if and only if $a\ge a_{\min}$. With $m=a/\gamma$,
$m_{\min}=a_{\min}/\{\beta(1-n)\}=n(2-n)/(1-n)^2$, which equals $F_\infty-1$ by
\cref{prop:th-spectrum}. Also $r m^2=r a^2/\gamma^2=\bar\Lambda n(2-n)/(1-n)^2$.

(c) For bins $I,J$, $\Cov\{N(I),N(J)\}$ integrates the covariance measure, which is
the inverse Fourier transform of the Bartlett spectrum. The variance-time curve
$\Var N(0,w]=\bar\Lambda w+2\int_0^w(w-u)c(u)\,\dd u$, the index of dispersion and the
expected periodogram are functionals of the same objects. Both $H$ and $C$ are
stationary Poisson cluster processes with almost surely finite clusters, so they are
mixing \citep[Ch.~12]{daley2008}. Their binned count sequences are therefore stationary
and ergodic, with finite second moments. By Birkhoff's theorem, empirical means and
autocovariances converge almost surely to the common population values. The
continuous mapping theorem concludes the proof. \qed

\subsection{Proof of \texorpdfstring{\cref{prop:th-cm}}{Proposition (CM kernels)}}

\emph{Step 1: finite mixtures.} Let $\varphi(t)=\sum_{i=1}^Kc_ie^{-b_it}$ with
$0<b_1<\dots<b_K$, $c_i>0$ and $\sum_ic_i/b_i=n<1$. Its Laplace transform
$\tilde\varphi(z)=\sum_ic_i/(z+b_i)$ is strictly decreasing on each interval between
consecutive poles. On $(-b_1,\infty)$ it falls from $+\infty$ to $0$ and equals $n<1$
at $z=0$. Hence $\tilde\varphi=1$ at exactly one point $z_1\in(-b_1,0)$. On each
$(-b_{i+1},-b_i)$ it falls from $+\infty$ to $-\infty$, which gives one more root
$z_{i+1}$.

Write $1-\tilde\varphi=P_1/P_0$ with $P_0(z)=\prod_i(z+b_i)$, so that $P_1$ is monic of
degree $K$. These $K$ roots are then all the roots of $P_1$; they are real, negative and simple.
Moreover $P_0-P_1$ has degree $K-1$, so
\[
\tilde R=\frac{\tilde\varphi}{1-\tilde\varphi}=\frac{P_0-P_1}{P_1}
=\sum_{l=1}^K\frac{\upsilon_l}{z-z_l},\qquad
\upsilon_l=\frac{1}{-\tilde\varphi'(z_l)}>0,
\]
since $\tilde\varphi(z_l)=1$ and $\tilde\varphi'<0$. The function
$\sum_l\upsilon_le^{z_lt}$ and the resolvent $R=\sum_j\varphi^{*j}\in L^1$ have the same
Laplace transform on $z\ge0$. Hence $R(t)=\sum_l\upsilon_le^{z_lt}$, a positive combination of
decaying exponentials, which is completely monotone. Also
$R(0+)=\lim_{z\to\infty}z\tilde R(z)=\varphi(0+)$.

\emph{Step 2: general completely monotone kernels.} By Bernstein's theorem,
$\varphi(t)=n\int_{(0,\infty)}e^{-bt}\sigma(\dd b)$ with $\int\sigma(\dd b)/b=1$. Fix
$\iota\ge1$ and partition $(2^{-\iota},2^\iota]$ into cells of width $2^{-\iota}$. Let
\[
\varphi_\iota(t)=n\sum_I\sigma(I)e^{-b_I^+t},
\]
where $b_I^+$ is the right end of cell $I$. Each $\varphi_\iota$ is a finite exponential
mixture. Refining a cell or extending the range can only increase $\varphi_\iota$, so
$\varphi_\iota\le\varphi_{\iota+1}\le\varphi$, and $\varphi_\iota\uparrow\varphi$ pointwise by
monotone convergence. Hence $\int\varphi_\iota\le n<1$ and, by Step 1, each
$R_\iota=\sum_k\varphi_\iota^{*k}$ is completely monotone.

Convolution powers are monotone in the kernel, so $R_\iota\uparrow R$ pointwise by monotone
convergence. The limit $R$ is nonincreasing and integrable, with $\int R=n/(1-n)$, so it
is finite on $(0,\infty)$. A pointwise limit of completely monotone functions is
completely monotone \citep[Cor.~1.6]{schilling2012}. Moreover
\[
R(0+)=\sup_{\iota}R_\iota(0+)=\sup_{\iota}\varphi_\iota(0+)\le\varphi(0+)<\infty .
\]

\emph{Step 3: covariance.} By \cref{prop:th-spectrum},
$2\pi f_H=\bar\Lambda|1+\hat R|^2$, so for $u\ne0$
\[
c_H(u)=\bar\Lambda\Big[R(|u|)+\int_0^\infty R(v)R(v+|u|)\,\dd v\Big].
\]
Write $R(t)=\int e^{-bt}\varpi(\dd b)$. By Tonelli,
$\int_0^\infty R(v)R(v+u)\,\dd v=\iint e^{-b'u}(b+b')^{-1}\varpi(\dd b)\varpi(\dd b')$.
Hence $c_H(u)=\int e^{-b|u|}\mathcal K(\dd b)$ with
\[
\mathcal K(\dd b)=\bar\Lambda\Big[1+\int\frac{\varpi(\dd b')}{b+b'}\Big]\varpi(\dd b)\ge0 .
\]
If $\varpi$ has $K$ atoms, so does $\mathcal K$. Also
$Z=\int \mathcal K(\dd b)/b=\tfrac12\int_{\mathbb R}c_H=\tfrac12\{2\pi f_H(0)-\bar\Lambda\}=\bar\Lambda(F_\infty-1)/2$,
and $\int \mathcal K(\dd b)=c_H(0)<\infty$ by Step 2.

\emph{Step 4: a Cox process with the same second-order structure.} Each shot integrates
to $m$, so $\E\xi=\nu+r m=\bar\Lambda$. By Campbell's theorem for the marked Poisson
process $\{(c_\ell,\Gamma_\ell)\}$, for $u\ge0$,
\[
\Cov\{\xi(t),\xi(t+u)\}=r\,\E\int_{-\infty}^tm^2\Gamma^2e^{-\Gamma(t-c)}e^{-\Gamma(t+u-c)}\,\dd c
=\frac{r m^2}{2}\int b\,e^{-bu}\,\pi(\dd b)=\frac{r m^2}{2Z}\int e^{-bu}\mathcal K(\dd b),
\]
which equals $c_H(u)$ for $r=2Z/m^2$. Then
$\nu=\bar\Lambda-2Z/m=\bar\Lambda\{1-(F_\infty-1)/m\}\ge0$ if and only if
$m\ge F_\infty-1$. Also $\E\xi^2<\infty$ because $c_H(0)<\infty$. By
\cref{lem:th-cox2}(a), the Cox process has mean $\bar\Lambda$ and covariance density
$c_H$, so it is second-order equivalent to $H$.

\emph{Step 5.} $2\pi f_H(\omega)=\bar\Lambda+\int2b(b^2+\omega^2)^{-1}\mathcal K(\dd b)\ge\bar\Lambda$.

\emph{Gamma kernel (\cref{rem:th-noncm}).} For $g(t)=\beta^2te^{-\beta t}$, put
$x=\omega^2/\beta^2$. Then
\[
|1-\hat\varphi|^2=\frac{(1-n-x)^2+4x}{(1+x)^2}.
\]
Its derivative in $x$ has the sign of $3-n-x$, so it is maximised at $x=3-n$, where it
equals $4/(4-n)$. Hence $\min_\omega2\pi f_H/\bar\Lambda=1-n/4$. \qed

\subsection{Proof of \texorpdfstring{\cref{prop:th-mv}}{Proposition (multivariate)}}

(a) This is \cref{lem:th-cox2}(a).

(b) Row $k$ of $\Imat-\hat\Phik$ is $\bm e_k^{\top}$, so
$\bm e_k^{\top}=\bm e_k^{\top}\bm G$. Hence $(\bm G\bm\Delta\bm G^*)_{kk}=\bar\Lambda_k$ and
$(\bm G\bm\Delta\bm G^*)_{jk}=G_{jk}\bar\Lambda_k$, where
$\bm\Delta=\operatorname{diag}(\bar{\bm\Lambda})$. Thus $E_{H,kk}\equiv0$ and
$E_{H,jk}=G_{jk}\bar\Lambda_k$. At $\omega=0$,
$\bm G(0)=\Mmat=\sum_m\Amat^m\ge\Amat$ entrywise, so $G_{jk}(0)\ge\alpha_{jk}>0$.
Since $\hat\Phik$ is continuous, $G_{jk}\neq0$ on a neighbourhood $U$ of $0$. On $U$,
the principal $2\times2$ submatrix of $\bm E_H$ indexed by $\{k,j\}$ has determinant
$-|E_{H,jk}|^2<0$, so $\bm E_H$ is not positive semidefinite there. By (a), no Cox
process has this spectrum.

(c) Campbell's theorem gives
$\Cov\{\xi_k(t+u),\xi_j(t)\}=r\E[a_ka_j]\int s_k(t+u-c)s_j(t-c)\,\dd c$. \qed

\subsection{Proof of \texorpdfstring{\cref{thm:spurious}}{Theorem (spurious contagion)}}

(a) By \cref{lem:th-cox2}, $2\pi f_N(\omega)=\bar\Lambda+2\sigma^2\gamma/(\gamma^2+\omega^2)$.
An exponential Hawkes process has
$2\pi f(\omega)=\bar\Lambda_H+\bar\Lambda_H\beta^2n(2-n)/\{\beta^2(1-n)^2+\omega^2\}$.
As in the proof of \cref{thm:equiv}(a), second-order equivalence holds if and only if
$\bar\Lambda_H=\bar\Lambda$, $\beta(1-n)=\gamma$ and
$\bar\Lambda\beta^2n(2-n)=2\sigma^2\gamma$. Using
$\beta^2n(2-n)=\gamma^2\{(1-n)^{-2}-1\}$, the last condition becomes
\[
(1-n)^{-2}=1+2\sigma^2/(\bar\Lambda\gamma)=1+\kappa ,
\]
since $\int c_\xi=2\sigma^2/\gamma$. This has the unique solution $n^*\in(0,1)$, and
then $\beta^*=\gamma/(1-n^*)$ and $\mu^*=\bar\Lambda(1-n^*)$. For the shot-noise process,
$\sigma^2=r a^2/(2\gamma)$. For the two-state chain,
$\Pr(X_u=1\mid X_0=1)=\pi_1+\pi_0e^{-(q_{01}+q_{10})u}$, which gives
$c_\xi(u)=(\lambda_1-\lambda_0)^2\pi_0\pi_1e^{-(q_{01}+q_{10})|u|}$.

(b) Let $\bm m(\vartheta)$ denote the population counterpart of $\hat{\bm m}_T$. Since $N$
is stationary and ergodic with finite second moments, $\hat{\bm m}_T\to\bm m(\vartheta^*)$
almost surely under $N$, by Birkhoff's theorem and (a). Under $H(\vartheta^*)$, which is
ergodic (proof of \cref{thm:equiv}(c)), also $\hat{\bm m}_T\to\bm m(\vartheta^*)$, and by
assumption $G(\hat{\bm m}_T)\to\vartheta^*$. Continuity of $G$ then forces
$G(\bm m(\vartheta^*))=\vartheta^*$. Hence, under $N$,
$\hat\vartheta_T=G(\hat{\bm m}_T)\to G(\bm m(\vartheta^*))=\vartheta^*$.

(c) $\hat c_\xi$ is continuous because $c_\xi\in L^1$, so
$F_\infty(N)=1+\hat c_\xi(0)/\bar\Lambda=1+\kappa$. For a Hawkes process,
$F_\infty=|1-\hat\varphi(0)|^{-2}=(1-n)^{-2}$, so $n_F=n$.

(d) Fix $\beta$ and set $X(t)=\int_{(-\infty,t)}g_\beta(t-s)\,\dN(s)\ge0$, a predictable
stationary process with $\E X=\bar\Lambda$. The second factorial moment measure of $N$
has density $\bar\Lambda^2+c(t-s)$ off the diagonal, so
\begin{equation}\label{eq:th-EXdN}
\E\int_0^1X(t)\,\dN(t)=\int_0^1\int_{-\infty}^{t}g_\beta(t-s)\{\bar\Lambda^2+c(t-s)\}\,\dd s\,\dd t
=\bar\Lambda^2+I_\beta,
\qquad I_\beta=\int_0^\infty g_\beta c .
\end{equation}
Since $\log\mu\le\log\lambda_\vartheta\le\log\mu+nX/\mu$, the function $L$ is finite.

For $n=0$, $L(\mu,0,\beta)=\bar\Lambda\log\mu-\mu$, which is uniquely maximised at
$\mu=\bar\Lambda$ with value $\bar\Lambda\log\bar\Lambda-\bar\Lambda$. Now fix
$\mu=\bar\Lambda$ and $\beta=\beta_0$. By concavity, the difference quotients
$D_n=\{\log(\bar\Lambda+nX)-\log\bar\Lambda\}/n\ge0$ increase to $X/\bar\Lambda$ as
$n\downarrow0$. By monotone convergence and \eqref{eq:th-EXdN},
\[
\frac{L(\bar\Lambda,n,\beta_0)-L(\bar\Lambda,0,\beta_0)}{n}
=\E\int_0^1D_n\,\dN-\bar\Lambda\ \longrightarrow\ \frac{I_{\beta_0}}{\bar\Lambda}>0 .
\]
Hence $L(\bar\Lambda,n,\beta_0)>\bar\Lambda\log\bar\Lambda-\bar\Lambda$ for all small
$n>0$. For $C$, $c(u)=\{r a^2/(2\gamma)\}e^{-\gamma|u|}$, and integrating against
$\beta e^{-\beta u}$ gives the stated $I_\beta$. \qed

\subsection{Proof of \texorpdfstring{\cref{cor:equiv-id}}{Corollary (identification)}}

(i) Under \cref{ass:th-model} the Volterra resolvent of $\Phik$ on $[0,T]$ is
integrable, so $\E\bm N[0,T]<\infty$, and the intensity has the form of
\cref{lem:th-janossy} with
\[
h_k(t;x)=m_k(t)+\sum_j\sum_{s\in x_j}\varphi_{kj}(t-s),\qquad h_k\ge\inf\mu_k>0 .
\]
Take two parameter values with the same law. Taking $x=\emptyset$ in
\cref{lem:th-janossy} gives $m^{(1)}_k=m^{(2)}_k$ almost everywhere on $[0,T]$. Taking
$x=\{(s,j)\}$ then gives $\varphi^{(1)}_{kj}(t-s)=\varphi^{(2)}_{kj}(t-s)$ for almost
every $0<s<t<T$. By Fubini, $\varphi^{(1)}_{kj}=\varphi^{(2)}_{kj}$ almost everywhere
on $(0,T)$.

For the sufficient condition, $m_k=\mu_k$ on $(0,\tau_1)$. This determines $\mu_k$ when
$\mu_k$ is constant, and on all of $[0,T]$ by analytic continuation when it is analytic.
Then $\sum_eS_e\beta_k\psi_k(t-\tau_e)=m_k(t)-\mu_k(t)$ is known. On
$[\tau_1,\tau_2)$ only the first shock is active, so this determines the product
$\beta_kS_1\psi_k$ there. Since $S_1>0$ is observed, it determines $\beta_k\psi_k$ on
$[0,\tau_2-\tau_1)$. By the condition on the scaled family, $\beta_k\psi_k$ is then
determined on $(0,\infty)$. Then $\int\psi_k=1$ gives
$\beta_k=\int\beta_k\psi_k$, and hence $\psi_k$. Finally, if $\alpha_{kj}>0$ and $g_{kj}$ does not
vanish on $(0,T)$, then $\varphi_{kj}\ne0$ on a set of positive measure, so the law
differs from that with $\alpha_{kj}=0$.

(ii) By \cref{lem:th-cluster}(a) and \cref{prop:stationary}(c), the cumulant rates of $H$
are $\mu\E \mathcal S^j$. Normalise $\mu=1$ (all $\chi_j$ scale linearly in $\mu$) and write
$u=1-n$. Then
\[
\chi_1=u^{-1},\quad \chi_2=u^{-3},\quad \chi_3=(3-2u)u^{-5},\quad
\chi_4=(15-20u+6u^2)u^{-7}.
\]
The cumulant rates of $C$ are $\nu+r\E[\Pois(m)^j]$, with
$\E\Pois(m)^j=m$, $m+m^2$, $m+3m^2+m^3$ and $m+7m^2+6m^3+m^4$ for $j=1,\dots,4$.
Matching $j=1,2,3$ gives $r m^2=\chi_2-\chi_1$ and
$r(3m^2+m^3)=\chi_3-\chi_1$, hence
\[
m+3=\frac{\chi_3-\chi_1}{\chi_2-\chi_1}
=\frac{3+u+u^2+u^3}{u^2(1+u)},\qquad
m=\frac{(1-u)(3+4u+2u^2)}{u^2(1+u)}>0 .
\]
This determines $(\nu,r,m)$ uniquely, with $\gamma$ free. Moreover
\[
m-m_{\min}=\frac{(1-u)(2+2u+u^2)}{u^2(1+u)}>0 ,
\]
so with $\gamma=\beta(1-n)$ this is a member of the family of \cref{thm:equiv}. Since
$7m^2+6m^3+m^4=m^2\{(m+3)^2-2\}$,
\[
\chi_4^C=\chi_1+\frac{(\chi_3-\chi_1)^2}{\chi_2-\chi_1}-2(\chi_2-\chi_1),
\qquad
\chi_4^H-\chi_4^C=\frac{\mathcal D}{\chi_2-\chi_1},
\]
where
\[
\mathcal D=(\chi_4-\chi_1)(\chi_2-\chi_1)+2(\chi_2-\chi_1)^2-(\chi_3-\chi_1)^2 .
\]
Direct expansion gives
\[
u^{10}\mathcal D=6-8u-13u^2+20u^3+2u^4-4u^5-5u^6+2u^8=(1-u)^4\,Q(u).
\]
Using $\chi_2-\chi_1=n(2-n)(1-n)^{-3}$ yields the stated difference, which is
positive for $n\in(0,1)$ since $Q>0$ on $(0,1)$. If $\Pr_H=\Pr_C$, all cumulant rates
would coincide, which is impossible.

A second proof uses \cref{lem:th-cluster}(b). For the Borel law,
$\Pr(\mathcal S=2)\Pr(\mathcal S=4)/\Pr(\mathcal S=3)^2=32/27$. For the cluster law of $C$ the same ratio is $3/4$,
since the $\nu$-component only affects $k=1$.

(iii) In the cluster representation, each stream-$j$ incident has $\Pois(\alpha_{kj})$
stream-$k$ offspring, independently of everything else. By stationarity, the
parent--offspring pairs of type $(j,k)$ therefore occur at rate $\alpha_{kj}\bar\Lambda_j$
per unit time, and each is recorded with probability $q$. \qed

\subsection{Proof of \texorpdfstring{\cref{thm:thinning}}{Theorem (thinning)} and \texorpdfstring{\cref{cor:th-darkfig}}{Corollary (dark-figure correction)}}

(a) This follows from \cref{lem:th-cox2}(b) with $M=1$ and \eqref{eq:th-fH}.

(b) By the argument of \cref{thm:equiv}(a), $H(\mu',n',\beta')$ is second-order
equivalent to $N^p$ if and only if $\bar\Lambda'=p\bar\Lambda$,
$\beta'(1-n')=\beta(1-n)=:b$ and $\beta'^2n'(2-n')=p\beta^2n(2-n)$. Put
$F(x)=(1-x)^{-2}$. Then $\beta^2n(2-n)=b^2\{F(n)-1\}$ and similarly for $n'$, so the
last condition reads $F(n')-1=p\{F(n)-1\}$. This determines $n'$ uniquely, and then
$\beta'^2=b^2F(n')=\beta^2(1-n)^2+p\beta^2n(2-n)$ and $\mu'=\bar\Lambda'(1-n')$.

(c) The function $F$ is a strictly increasing, strictly convex bijection
$[0,1)\to[1,\infty)$ with $F(0)=1$, and
\[
h_p=F^{-1}\circ\{1+p(F-1)\}.
\]
Hence $h_p$ is an increasing bijection. Next, $h_p(n)<n$ if and only if
$(1-p)\{F(n)-1\}>0$. Also $h_p(n)>pn$ if and only if
$F(pn)<(1-p)F(0)+pF(n)$, which is strict convexity. The expansion follows from
$F(n)=1+2n+O(n^2)$. Finally,
\[
\beta'=b\,F(n')^{1/2}=\beta\{F(n')/F(n)\}^{1/2}\in[\beta(1-n),\beta],
\]
and solving $F(n)=1+\{F(n')-1\}/p$ gives \eqref{eq:th-hinv}.

(d) $N^p$ is a Poisson cluster process with $r=\mu$ and cluster size
$\mathcal S_p\mid \mathcal S\sim\Bin(\mathcal S,p)$, where $\mathcal S\sim$ Borel$(n)$, so $\E e^{u \mathcal S_p}<\infty$ for
small $u>0$. Suppose $N^p$ were a stationary Hawkes process with constant
background $\mu''\ge0$ and some kernel $\varphi''\ge0$ of mass $n''$.

First suppose $\mu''=0$. Taking expectations of the intensity gives
$p\bar\Lambda=n''p\bar\Lambda$, so $n''=1$. The process $N^p$ has finite second
moments, so $\dN^p-\lambda''\dd t$ is a martingale measure with spectral density
$p\bar\Lambda/(2\pi)$. Since $\dN^p-\lambda''\dd t$ is the output of the filter
$\delta_0-\varphi''$ applied to $\dN^p$, its spectral measure is
$|1-\hat\varphi''(\omega)|^2f_p(\omega)\,\dd\omega$. Hence
$|1-\hat\varphi''(\omega)|^2f_p(\omega)=p\bar\Lambda/(2\pi)$ for almost every $\omega$.
As $\omega\to0$, $\hat\varphi''(\omega)\to n''=1$ while $f_p$ stays bounded by (a),
since $F'_\infty<\infty$. The left side therefore tends to $0$, a contradiction.

So $\mu''>0$. Then \cref{prop:stationary}(b) gives $n''<1$, and $N^p$ is a Poisson
cluster process with Borel$(n'')$ clusters. By
\cref{lem:th-cluster}(b), $\mu\Pr(\mathcal S_p=k)=\mu''\Pr(\text{Borel}(n'')=k)$ for $k\ge1$.

For any Borel law with $n''>0$,
\[
\frac{\Pr(1)\Pr(3)}{\Pr(2)^2}=\frac{e^{-n''}\tfrac32n''^2e^{-3n''}}{(n''e^{-2n''})^2}=\frac32 .
\]
If $n''=0$, then $\Pr(2)=0<\Pr(\mathcal S_p=2)$, which is already a contradiction. For $\mathcal S_p$,
write the Borel generating function as $G(z)=\Theta(\zeta z)/n$, where $\zeta=ne^{-n}$ and
$\Theta(y)=\sum_{k\ge1}k^{k-1}y^k/k!$ is the tree function, which satisfies $\Theta=ye^\Theta$. Then
$\Pr(\mathcal S_p=k)=p^kG^{(k)}(1-p)/k!=(p\zeta)^k\Theta^{(k)}(y)/(n\,k!)$ with
$y=\zeta(1-p)\in(0,e^{-1})$. Differentiating $\Theta=ye^\Theta$ and writing $v=\Theta(y)\in(0,1)$,
\[
\Theta'=\frac{e^v}{1-v},\qquad \Theta''=\frac{e^{2v}(2-v)}{(1-v)^3},\qquad
\Theta'''=\frac{e^{3v}(2v^2-8v+9)}{(1-v)^5}.
\]
Therefore
\[
\frac{\Pr(\mathcal S_p=1)\Pr(\mathcal S_p=3)}{\Pr(\mathcal S_p=2)^2}=\frac23\,\frac{\Theta'\Theta'''}{(\Theta'')^2}
=\frac23\,\frac{2v^2-8v+9}{(2-v)^2}
=\frac32+\frac{v(4-v)}{6(2-v)^2}>\frac32 ,
\]
a contradiction. The statement about moment-based estimators follows as in
\cref{thm:spurious}(b), since $N^p$ is a Poisson cluster process (hence ergodic) and is
second-order equivalent to $H(\mu',n',\beta')$.

\emph{Proof of \cref{cor:th-darkfig}.} Write $h^{-1}(n',p)=1-\mathcal F^{-1/2}$ with
$\mathcal F=1+\{(1-n')^{-2}-1\}/p=(1-n)^{-2}$. Then
\[
\frac{\partial h^{-1}}{\partial n'}=\frac{\mathcal F^{-3/2}}{p\,(1-n')^3}=\frac{(1-n)^3}{p\,(1-n')^3},
\qquad
\frac{\partial h^{-1}}{\partial p}=-\frac{\mathcal F^{-3/2}\{(1-n')^{-2}-1\}}{2p^2}
=-\frac{n(1-n)(2-n)}{2p},
\]
using $(1-n')^{-2}-1=p\{(1-n)^{-2}-1\}$. The map $h^{-1}$ is $C^1$ on
$(0,1)\times(0,1]$. Consistency follows from the continuous mapping theorem, and
\eqref{eq:th-delta} from the delta method and independence. \qed

\subsection{Proof of \texorpdfstring{\cref{prop:th-hetthin}}{Proposition (heterogeneous recording)}}

Marks are unpredictable with law $\pi$, so incidents with mark $x$ have intensity
$\pi_x\{\mu+\sum_in_{X_i}g(t-t_i)\}$. The marked process is therefore a $K$-variate
Hawkes process with background $\mu\bm\pi$, branching matrix $\Amat=\bm\pi\bm n^{\top}$
and common kernel $g$. Since $\Amat$ has rank one, $\spr(\Amat)=\bm n^{\top}\bm\pi=\bar n<1$.
By the Sherman--Morrison formula,
$(\Imat-\Amat)^{-1}=\Imat+\bm\pi\bm n^{\top}/(1-\bar n)$, so
$\bar{\bm\Lambda}=\mu\bm\pi/(1-\bar n)=\bar\Lambda\bm\pi$. Similarly, with
$\hat\Phik=\hat g\,\bm\pi\bm n^{\top}$,
\[
\bm G=\Imat+\mathsf H\bm\pi\bm n^{\top},\qquad \mathsf H=\hat g/(1-\bar n\hat g).
\]
The recorded process is $\mathbf 1^{\top}\bm N^{\bm P}$ with $\bm P=\operatorname{diag}(p_x)$.
By \cref{lem:th-cox2}(b) and \cref{prop:th-spectrum}, with
$\bm\Delta=\bar\Lambda\operatorname{diag}(\bm\pi)$ and $\bm p=(p_x)$,
\[
2\pi f_{\mathrm{rec}}=\bm p^{\top}\bar{\bm\Lambda}+\bm p^{\top}(\bm G\bm\Delta\bm G^*-\bm\Delta)\bm p .
\]
Since
$\bm G\bm\Delta\bm G^*-\bm\Delta=\mathsf H\bm\pi\bm n^{\top}\bm\Delta+\bar{\mathsf H}\bm\Delta\bm n\bm\pi^{\top}+|\mathsf H|^2(\bm n^{\top}\bm\Delta\bm n)\bm\pi\bm\pi^{\top}$,
and $\bm p^{\top}\bm\pi=\bar p$, $\bm n^{\top}\bm\Delta\bm p=\bar\Lambda\E[p_Xn_X]$,
$\bm n^{\top}\bm\Delta\bm n=\bar\Lambda\E[n_X^2]$,
\[
2\pi f_{\mathrm{rec}}=\bar p\bar\Lambda\big[1+2\operatorname{Re}(\mathsf H)\E[p_Xn_X]+|\mathsf H|^2\bar p\,\E[n_X^2]\big].
\]
For $g(t)=\beta e^{-\beta t}$, $\mathsf H=\beta/(b+i\omega)$ with $b=\beta(1-\bar n)$, so
$2\operatorname{Re}\mathsf H=2\beta b/(b^2+\omega^2)$ and
$|\mathsf H|^2=\beta^2/(b^2+\omega^2)$. This gives the stated spectrum. Matching with
an exponential Hawkes spectrum, as in \cref{thm:thinning}(b), gives $b''=b$ and
\[
(1-n'')^{-2}=1+\frac{2(1-\bar n)\E[p_Xn_X]+\bar p\,\E[n_X^2]}{(1-\bar n)^2}.
\]
Finally, $(1-n_{\mathrm{naive}})^{-2}=1+\{(1-n'')^{-2}-1\}/\bar p$. Substituting
$\E[p_Xn_X]=\bar p\bar n+\Cov(p_X,n_X)$, $\E[n_X^2]=\bar n^2+\Var(n_X)$ and
$1-(1-\bar n)^2=2\bar n(1-\bar n)+\bar n^2$ gives the stated expression. \qed

\subsection{Proof of \texorpdfstring{\cref{prop:granger}}{Proposition (Granger)}}

As in the proof of \cref{cor:equiv-id}(i), the intensity has the form of
\cref{lem:th-janossy} with
$h_k(t;x)=m_k(t)+\sum_l\sum_{s\in x_l}\varphi_{kl}(t-s)$.

If $\varphi_{kj}=0$ almost everywhere on $(0,T)$, then $h_k(t;x)$ depends on $x$ only
through $x_{-j}$, the configuration with stream $j$ removed. Hence $\lambda_k$ is
$\Hist^{V\setminus\{j\}}$-predictable.

Conversely, suppose $\lambda_k$ has a $\Hist^{V\setminus\{j\}}$-predictable version.
Predictable processes of a point-process history are measurable functionals of the
history strictly before $t$ \citep[Ch.~14]{daley2008}. So
$h_k(t;x)=\Psi(t,x_{-j})$ for $\Pr\otimes\dd t$-almost every $(t,x)$, and, because the
Janossy densities are positive, for almost every $(t,x)$ with respect to Lebesgue
measure times the Lebesgue--Poisson measure. The configurations $x=\emptyset$ and
$x=\{(s,j)\}$ have the same $x_{-j}=\emptyset$. Hence
\[
m_k(t)=\Psi(t,\emptyset)=m_k(t)+\varphi_{kj}(t-s)\qquad\text{for almost every }0<s<t<T,
\]
that is, $\varphi_{kj}=0$ almost everywhere on $(0,T)$. For exponential kernels,
$g_{kj}>0$ everywhere, so this is equivalent to $\alpha_{kj}=0$. \qed
